\appendix

\section{Pfaf\/f\/ians} \label{appendixA}

Let ${\mathcal A} =(a_{ij})_{1\le i,j\le2m}$ be a~skew-symmetric matrix.
Then the Pfaf\/f\/ian ${\rm Pf}\, {\mathcal A}$ is def\/ined by
\begin{gather*}
\det{\mathcal A}=({\rm Pf}\,{\mathcal A})^2,\qquad {\rm Pf}\,{\mathcal A}=a_{12}a_{34}\cdots a_{2m-1,2m}+\cdots.
\end{gather*}
Following~\cite{H1} we denote ${\rm Pf}\, {\mathcal A}$ by $(1,2,3,\dots,2m)$:
\begin{gather*}
{\rm Pf}\,{\mathcal A}=(1,2,3,\dots,2m).
\end{gather*}
It is directly def\/ined by
\begin{gather*}
(1,2,3,\dots,2m)=\sum{\rm sgn}(i_1,\dots,i_{2m})\cdot(i_1,i_2)(i_3,i_4)\cdots(i_{2m-1},i_{2m}
),\qquad (i,j)=a_{ij},
\end{gather*}
where the sum is over all permutations of $(1,\dots,2m)$ such that
\begin{gather*}
i_1<i_3<\cdots<i_{2m-1},\qquad i_1<i_2,\dots,i_{2m-1}<i_{2m},
\end{gather*}
and ${\rm sgn}(i_1,\dots,i_{2m})$ is the signature of the permutation $(i_1,\dots,i_{2m})$.
The Pfaf\/f\/ian can be expanded as
\begin{gather*}
(1,2,3,\dots,2m)=\sum_{j=2}^{2m}(-1)^j(1,j)(2,3,\dots,\hat{j},\dots,2m).
\end{gather*}
For example, in the case of $m=2$,
\begin{gather*}
(1,2,3,4)=(1,2)(3,4)-(1,3)(2,4)+(1,4)(2,3).
\end{gather*}

\section{Sylvester's theorem for determinants and Pfaf\/f\/ians} \label{appendixB}

The following theorem is known as
Sylvester's theorem.
\begin{theorem}\label{theorem4}
Let $r\le m$, ${\mathcal A}=(a_{ij})_{1\le i,j\le m}$ and ${\mathcal A}_r=(a_{ij})_{1\le i,j\le r}$.
Set
\begin{gather*}
{\mathcal B}=(b_{ij})_{r+1\le i,j\le m},
\qquad
b_{ij}=\det\left(
\begin{matrix}a_{11}&\hdots&a_{1r}&a_{1j}
\\
\vdots&\ddots&\vdots&\vdots
\\
a_{r1}&\hdots&a_{rr}&a_{rj}
\\
a_{i1}&\hdots&a_{ir}&a_{ij}
\end{matrix}
\right).
\end{gather*}
Then we get
\begin{gather*}%\label{eq:90}
\det{\mathcal B}=(\det{\mathcal A}_r)^{m-r-1}\det{\mathcal A}.
\end{gather*}
\end{theorem}

  Let ${\mathcal A}=(a_{ij})_{1\le i,j\le2m}$ be a~skew-symmetric matrix and set $(i,j)=a_{ij}$.
For $r\le m$ let ${\mathcal P} =(p_{ij})_{2r+1\le i,j\le2m}$, $p_{ij}=(1,2,\dots,2r,i,j)$ and $I_r=\{1,2,\dots,2r\}$.
In general, for a~subset $I\subset\{1,2,\dots,2m\}$ we set ${\mathcal A}(I)=(a_{ij})_{i,j\in I}$ and for $i<j$, $k<l$
we denote by ${\mathcal A}_{kl}^{ij}$ be the square matrix of degree $2(m-1)$ which is obtained from ${\mathcal A}$ by
removing $i$-th and $j$-th rows, $k$-th and $l$-th columns.

\begin{theorem}[\cite{H2}]\label{theorem5} For $r\le m$ the following identity holds:
\begin{gather*}%\label{eq:91}
{\rm Pf}\,{\mathcal P}=({\rm Pf}\,{\mathcal A}(I_r))^{m-r-1}{\rm Pf}\,{\mathcal A}.
\end{gather*}
\end{theorem}

\section{The Pl\"{u}cker relation for Pfaf\/f\/ians}\label{appendixC}

 There exist analogues of the Pl\"{u}cker's
relations for Pfaf\/f\/ians~\cite{O1}.
They are given by
\begin{gather}
\sum_{l=1}^L(-1)^l(i_1,\dots,i_K,j_l)(j_1,\dots,\hat{j}_l,\dots,j_L)
\nonumber
\\
\qquad
{}+\sum_{k=1}
^K(-1)^k(i_1,\dots,\hat{i}_k,\dots,i_K)(j_1,\dots,j_L,i_k)=0,
\label{eq:75}
\end{gather}
where $K$ and $L$ are odd.
We understand that $(\varnothing)=1$.

For $n$ odd, taking $K=1$, $L=n$, $i_1=0$ and $j_1,\dots,j_n\neq0$ in~\eqref{eq:75}, we get
\begin{gather}
\sum_{l=1}^n(-1)^{l-1}(0,j_l)(j_1,\dots,\hat{j}_l,\dots,j_n)-(0,j_1,\dots,j_n)=0.
\label{eq:500}
\end{gather}
For $n$ even, setting $K=1$, $L=n-1$ and $i_1\neq0$ in~\eqref{eq:75}, we have
\begin{gather}
\sum_{l=1}^{n-1}(-1)^{l-1}(i_1,j_l)(j_1,\dots,\hat{j}_l,\dots,j_{n-1})-(i_1,j_1,\dots,j_{n-1})=0.
\label{eq:501}
\end{gather}
